\documentclass[12pt]{article}
\usepackage[T1]{fontenc}
\usepackage[utf8]{inputenc}
\usepackage{lmodern}
\usepackage{textcomp}
\usepackage{enumitem}
\usepackage{geometry}
\geometry{margin=1in}
\begin{document}
\begin{center}
Answer Key - Biology - Undergraduate Level Assessment 7 \
\small (Answers are ordered exactly as in the question paper.)
\end{center}

\section*{Multiple Choice}
\begin{enumerate}[label=\arabic*.]
\item \textbf{Answer:} A
\\ \textit{Options:} A) Levels of oxygen, acidity or alkalinity of the medium; B) Temperature changes, light exposure; C) Atmospheric pressure, noise levels; D) Nitrogen levels, radiation exposure; E) Soil texture, light wavelength
\\ \textit{Note:} Correct option: A - Levels of oxygen, acidity or alkalinity of the medium
\item \textbf{Answer:} C
\\ \textit{Options:} A) Lysosomes act as storage units for excess nutrients; B) Lysosomes are used for energy production; C) Lysosomes function in intra cellular digestion; D) Lysosomes synthesize proteins for the cell; E) Lysosomes regulate cell division
\\ \textit{Note:} Correct option: C - Lysosomes function in intra cellular digestion
\item \textbf{Answer:} A
\\ \textit{Options:} A) Only two DNA regions, V and D, combine in heavy chains; B) V, D, and JH regions are combined in light chains; C) V and JH regions are combined, while D regions are excluded from heavy chains; D) Five separate DNA regions are combined in heavy chains, adding extra complexity; E) D and JH regions are combined, while V regions are excluded from heavy chains
\\ \textit{Note:} Correct option: A - Only two DNA regions, V and D, combine in heavy chains
\item \textbf{Answer:} A
\\ \textit{Options:} A) 12.5 percent; B) 40 percent; C) 18.75 percent; D) 25 percent; E) 37.5 percent
\\ \textit{Note:} Correct option: A - 12.5 percent
\item \textbf{Answer:} E
\\ \textit{Options:} A) - .09176; B) .21185; C) .37748; D) .55962; E) .65361
\\ \textit{Note:} Correct option: E - .65361
\end{enumerate}

\section*{True / False}
\begin{enumerate}[label=\arabic*.]
\setcounter{enumi}{5}
\item \textbf{Answer:} False
\\ \textit{Options:} A) True; B) False
\\ \textit{Note:} LLM-generated false statement. Correct answer: 'directional selection'.
\item \textbf{Answer:} False
\\ \textit{Options:} A) True; B) False
\\ \textit{Note:} LLM-generated false statement. Correct answer: 'inheritance of both "fit" and "unfit" genes'.
\item \textbf{Answer:} False
\\ \textit{Options:} A) True; B) False
\\ \textit{Note:} LLM-generated false statement. Correct answer: 'Innate behavior'.
\item \textbf{Answer:} False
\\ \textit{Options:} A) True; B) False
\\ \textit{Note:} LLM-generated false statement. Correct answer: 'clathrin'.
\item \textbf{Answer:} True
\\ \textit{Options:} A) True; B) False
\\ \textit{Note:} LLM-generated true statement based on MCQ.
\end{enumerate}

\section*{Long Form Response}
\begin{enumerate}[label=\arabic*.]
\setcounter{enumi}{10}
\item \textbf{Answer:} The size of the appendix varies greatly among individuals.. Explain why this is the correct answer and provide supporting evidence. Reference key facts or concepts that support this choice.
\\ \textit{Note:} Long-form derived from MCQ; award credit for stating the correct idea and giving supporting reasoning.
\item \textbf{Answer:} A flock of birds can be considered an elementary form of society called a motion group.. Explain why this is the correct answer and provide supporting evidence. Reference key facts or concepts that support this choice.
\\ \textit{Note:} Long-form derived from MCQ; award credit for stating the correct idea and giving supporting reasoning.
\end{enumerate}

\vfill
\noindent\textit{End of Answer Key}
\end{document}